% Environmental Modelling & Software submission.
%
% elsarticle option notes:
%   preprint    single column, single spaced  <- used here
%   review      single column, double spaced, line numbers (add for review copy)
%   3p,twocolumn  the printed double-column layout -- NOT wanted here
%   authoryear  REQUIRED for (Author year) citations. elsarticle defaults to
%               numbered mode, so without it \citep renders as [5] no matter which
%               .bst is loaded, and elsarticle-harv silently produces a numbered
%               reference list.
%
% EMS does not require a specific reference format at submission; any consistent
% style is accepted and the journal style is applied at proof stage. The journal's
% published in-text form is author-date, so elsarticle-harv is used.
\documentclass[preprint,authoryear,11pt]{elsarticle}

\usepackage[T1]{fontenc}
\usepackage{amsmath}
\usepackage[strings]{underscore}
\usepackage{booktabs}
\usepackage{array}
\usepackage{siunitx}
% siunitx does not define imperial units. The groundwater and sewer models are in
% feet, so declare it rather than dropping those lengths to prose.
\DeclareSIUnit{\foot}{ft}
\usepackage{graphicx}
% elsarticle 3.5 has no native ORCID support; orcidlink renders the iD as a
% hyperlinked icon after the author name.
\usepackage{orcidlink}
% elsarticle loads natbib itself -- do not load it again or the options clash.
\setcitestyle{aysep={}}          % (Author year), no comma
\usepackage{xcolor}
\usepackage{hyperref}
% Long repository URLs in the software statement cannot break in a 5 in measure and
% overflow the margin; xurl allows a break at any character.
\usepackage{xurl}
\hypersetup{colorlinks=true, linkcolor={red!50!black},
  citecolor={blue!50!black}, urlcolor={blue!80!black}}
\newcolumntype{L}[1]{>{\raggedright\arraybackslash}p{#1}}

% Full-width figures are drawn at 190 mm, which is the journal's full-page artwork
% width and the width the figure scripts author to. The preprint text block is only
% 360 pt (5.0 in), so \includegraphics[width=\textwidth] would scale them to 67% and
% render 6.5 pt annotations at 4.3 pt -- artwork "where text is disproportionately
% small", which the guide lists among the things not to submit. Prose keeps the
% narrower measure, which is easier to read, and the figures overhang it
% symmetrically; on letter paper that still leaves 0.51 in of margin each side.
% The box is declared zero width so the overhang raises no overfull warning.
\newlength{\figfull}
\setlength{\figfull}{190mm}
\newcommand{\widefig}[1]{\centerline{\makebox[0pt][c]{%
  \includegraphics[width=\figfull]{#1}}}}

\journal{Environmental Modelling \& Software}

\begin{document}

\begin{frontmatter}

% Title note: "library interfaces" rather than "the Basic Model Interface".
% MODFLOW 6 and D-Flow FM implement BMI; SWMM does not, and is driven through the
% SWMM toolkit API, which supplies the same capabilities under different names.
% Naming BMI in the title would misstate one of the three. BMI is carried in the
% abstract and keywords instead, where it can be qualified.
\title{Dynamic coupling of coastal hydrodynamic, groundwater flow, and sewer
network models using library interfaces}

% Name spellings and affiliations below are taken from each author's ORCID record
% (pub.orcid.org/v3.0/<id>/person and /employments), not from correspondence, so
% they should be confirmed by each author before submission. Office assignments
% for the New York Water Science Center authors follow their ORCID employment
% city: Coram for Bayraktar, Troy for Jahn.
% Affiliation labels follow FIRST APPEARANCE in the author list, and the
% \affiliation declarations are given in that same order. elsarticle letters them in
% declaration order, so any other arrangement prints the superscripts out of
% sequence: with Prigiobbe added at inst5 in third position they read a, b, e, c, d,
% b. Inserting an author means renumbering from that author onward.
% The guide requires a corresponding author to be marked. The address below is a
% PLACEHOLDER and must be replaced before submission -- it is the only detail here
% not taken from an ORCID record.
\author[inst1]{Joseph D. Hughes\corref{cor1}\,\orcidlink{0000-0003-1311-2354}}
\cortext[cor1]{Corresponding author. \textit{Email address:} jdhughes@intera.com}
\author[inst2]{Liv Herdman\,\orcidlink{0000-0002-5444-6441}}
\author[inst3]{Valentina Prigiobbe\,\orcidlink{0000-0001-6603-4295}}
\author[inst4]{Martijn J. Russcher\,\orcidlink{0000-0001-8799-6514}}
\author[inst5]{Banu Bayraktar\,\orcidlink{0000-0003-3612-6767}}
\author[inst2]{Kalle L. Jahn\,\orcidlink{0000-0002-4976-0137}}

\affiliation[inst1]{organization={INTERA Incorporated}, addressline={927 W Belle Plaine Ave}, city={Chicago}, state={IL}, postcode={60613}, country={USA}}

\affiliation[inst2]{organization={U.S. Geological Survey, New York Water Science Center}, addressline={425 Jordan Road}, city={Troy}, state={NY}, postcode={12180}, country={USA}}

\affiliation[inst3]{organization={Department of Geosciences, University of Padua}, addressline={Via Gradenigo 6}, city={Padua}, postcode={35131}, country={Italy}}

\affiliation[inst4]{organization={Deltares}, addressline={Boussinesqweg 1}, city={Delft}, postcode={2629 HV}, country={The Netherlands}}

\affiliation[inst5]{organization={U.S. Geological Survey, New York Water Science Center}, addressline={2045 Route 112, Building 4}, city={Coram}, state={NY}, postcode={11727}, country={USA}}

% SUBMISSION REQUIREMENTS -- Environmental Modelling & Software (ISSN 1364-8152).
% All of the following was read off the guide for authors, 2026-08-06, superseding
% the earlier notes here that came from search summaries.
%
% IN THIS FILE
%   Abstract     150 words maximum. Currently 309, so it needs roughly halving.
%                Deferred until Results exists, because what survives the cut should
%                be decided by what the findings turn out to be. Must stand alone,
%                avoid references, and define any uncommon abbreviation.
%   Keywords     1 to 7; seven here. Avoid keywords joined by "and" or "of".
%   Sections     NUMBERED, and cross-referenced by number. This overrides the
%                unnumbered house style for this submission; back matter below is
%                excluded from the sequence.
%   Equations    Numbered in the order they are first referred to in the text.
%   Corresponding author must be marked. See \cortext below.
%
% SEPARATE FILES AT SUBMISSION -- none of these are part of this .tex
%   Highlights         REQUIRED. 3 to 5 bullets, 85 characters maximum each,
%                      including spaces. Filename must contain "highlights".
%   Graphical abstract REQUIRED. 531 x 1328 px (h x w) or proportionally more,
%                      legible at 5 x 13 cm. TIFF, EPS, PDF or MS Office.
%   Figures            Separate files named Figure_1, Figure_2, ... Vector art as
%                      EPS or PDF with fonts embedded, which is what the scripts in
%                      scripts/ already produce.
%   Declarations       Competing interests, via Elsevier's declarations tool, as a
%                      .doc/.docx upload.
%
% STILL OUTSTANDING -- see the back matter for the sections themselves
%   Software and data availability, CRediT, and the generative-AI declaration are
%   drafted below but need author input before submission.
%   Research data: this journal is Option C, so data must be deposited in a
%   repository and cited, or a statement given explaining why it cannot be shared.
%   "Contact the author" is explicitly not acceptable for software or data access.
%
% NO PAGE CHARGES. The guide sets no page limit and no excess-page fee; the only
% fee is the optional open-access APC. Read in full, so this is now a confirmed
% absence rather than a failure to find one.
%
%   https://www.sciencedirect.com/journal/environmental-modelling-and-software/publish/guide-for-authors
\begin{abstract}
Hydrodynamic and hydrologic processes that span the coast, the subsurface, and engineered infrastructure are simulated by codes developed independently of one another, each with its own numerical formulation, grid, model time step, and control interface. An approach is described in which three such codes---a coastal hydrodynamic model, a groundwater flow model, and a sewer network model---are loaded as libraries into a single process and exchange state in memory at a fixed interval throughout a simulation, so that a self-consistent coupled solution is obtained from one simulation. None of the codes is modified: each is used as distributed and driven through its published interface, so the coupling does not fork a source tree or reimplement one process inside another, and it survives a new release of any component. Two of the three implement the Basic Model Interface and the third does not, which is shown to be immaterial: what a simulator must supply is the ability to be initialized, advanced and finalized as a shared library, and to have named state read and written, whatever those calls are named. Exchange formulations are given for each pair of models. Particular attention is given to returning the sewer model outfall discharge to the hydrodynamic model as a sewer discharge tracer source, which places the strongest requirement on the interface: the transfer has to occur between the external forcing update and the flow computation, a window the combined time-stepping entry point does not expose, and the value applied over a model time step is formed from the cumulative outfall volume so that mass is conserved exactly rather than sampled. The approach is demonstrated on a combined sewer network on eastern Long Island, New York, where groundwater infiltration to the sewer is modulated by the tide, and the sensitivity of the coupled solution to the exchange interval is evaluated across three hydrodynamic grid resolutions.
\end{abstract}

\begin{keyword}
model coupling \sep
Basic Model Interface \sep
model interoperability \sep
groundwater--surface water interaction \sep
D-Flow FM \sep
MODFLOW 6 \sep
SWMM
\end{keyword}

\end{frontmatter}

\section{Introduction}

Processes that cross the boundary between the coastal ocean, the subsurface, and engineered infrastructure are rarely represented by a single code. Each domain has a mature simulator of its own, developed independently, with its own governing equations, spatial discretization, model time step, and means of external control. Representing the coupled process therefore means making those codes exchange state while they run.

Coupling of this kind has generally been achieved in one of two ways. The first adds the second process to an existing code, which yields a tightly integrated solution but confines it to the domains the host code implements; the coupled surface water routing added to MODFLOW is an example \citep{hughes2014swr}. The second couples separate codes externally, and has been applied to urban drainage and groundwater through SWMM and MODFLOW \citep{zhang2020swmmmodflow, li2022urbansaturated, tajari2026compound}, to large-scale hydrology and groundwater \citep{guillaumot2022cwatm}, and to the coastal zone by coupling a hydrodynamic model to a groundwater model \citep{jin2025gwocean}.

Standardized control interfaces make the second route practical without the cost that has historically accompanied it. Simulators that expose a library interface, such as the Basic Model Interface \citep{hutton2020bmi} or an application programming interface built for the same purpose \citep{hughes2022modflowapi}, can be loaded into a single process and advanced under external control, with state passed as arrays in memory. The exchange interval is then set by the numerics rather than by any intermediate file format, and a self-consistent coupled solution is obtained from one simulation.

A simulator need not implement the Basic Model Interface as such to take part. The requirement is functional rather than nominal: the code must be callable as a shared library and must provide four capabilities---initialize a simulation from its own input, advance it by a time step, finalize it, and get and set named state. The Basic Model Interface is a codification of that set, and where a simulator implements it the coupling is written directly against it. Where a simulator does not, an interface offering the same four capabilities under different names serves equally well. Of the three codes coupled here, two implement the Basic Model Interface and the third does not, and the distinction turns out to matter only to the names of the calls. What does matter, and is returned to below, is whether the interface also permits access part way through a time step; that is a stronger requirement than the four capabilities above, and it is not part of the standard.

The decisive property is that none of the coupled codes is modified. Each is used as distributed: the coupling lives entirely in the controlling program, which advances each simulator through its published interface and moves state between them. No source is forked, no process is reimplemented inside a host code, and each simulator continues to be developed, tested and released by the group that maintains it. A coupling built this way survives a new release of any component, and the same controlling program can be repointed at a different combination of simulators. \cite{jin2025gwocean} exploit the same property in coupling a hydrodynamic model to MODFLOW~6, adjusting the general-head boundary through the interface rather than by altering the code.

What such an interface exposes, however, is not always sufficient on its own. A quantity may be readable but not writable in a way the solution uses, and where within the model time step a value is written can determine whether it survives to be used at all. The Basic Model Interface defines initialization, advancement in time, and finalization; it does not define access to model state part way through a time step. That omission is not incidental, and it has been encountered before: the MODFLOW application programming interface extends the Basic Model Interface precisely because the ability to modify variables within a time step falls outside it \citep{hughes2022modflowapi}. Any coupling that must inject state between two operations of a single time step therefore depends on something beyond the standard, and a simulator exposing only the standard interface cannot support that class of coupling at all. Both simulators driven here provide such an extension, by different routes, and the coupling described below depends on it.

The purpose of this paper is to describe an approach for coupling a coastal hydrodynamic model, a groundwater flow model, and a sewer network model in memory through their library interfaces, and to set out the exchange formulations between each pair. Returning the sewer model outfall discharge to the hydrodynamic model as a sewer discharge tracer source is treated in most detail, because it places the strongest requirement on the interface: the transfer must occur between the external forcing update and the flow computation, a window the combined time-stepping entry point such interfaces usually present does not expose.

The approach is demonstrated on a problem that exercises all three couplings at once. Sewers in low-lying coastal communities are subject to groundwater infiltration that varies with the tide: where the water table stands above the pipe invert, groundwater enters the sewer through defects in the pipe wall and at joints, adding to the volume the network conveys. The water table is itself tidally forced, so the infiltration rate is modulated at tidal frequencies by a process external to the sewer network. That interaction, and its consequences for flooding and for water quality, has been studied in coastal urban settings through field measurement, statistical and machine-learning models, and scenarios of sea level rise \citep{liu2018gwsewer, su2019infiltration, liu2021predicting, su2022slr}.

The demonstration uses the sewer network at Greenport, New York, on the north fork of eastern Long Island, which discharges to Greenport Harbor and the surrounding waters of Long Island Sound and Peconic Bay. The network is a combined system carrying both stormwater runoff and dry weather flow. It is used here to evaluate the coupling rather than to characterize the site, and in particular to establish how frequently the exchange must occur for the coupled solution to be insensitive to the exchange interval.

\section{Model components}

% TODO: the D-Flow FM Long Island Sound grid has no citable reference. Searches of
% Herdman's published works, USGS publications and Crossref turned up nothing that
% documents it, and the nearest Long Island compound-flooding paper
% (Glas et al., NHESS 2026) is statistical rather than hydrodynamic. If the grid is
% documented in a data release or report, cite it here; otherwise it needs a
% sentence describing its provenance.
Three simulators are coupled. Coastal hydrodynamics and the transport of a conservative tracer of sewer discharge are simulated with D-Flow Flexible Mesh (D-Flow FM), which solves the depth-averaged shallow water equations on an unstructured grid using the scheme of \cite{kernkamp2011unstructured} and is documented in \cite{deltares2026dflowfm}. Groundwater flow is simulated with MODFLOW~6 \citep{langevin2017modflow6}, using the Groundwater Flow (GWF) and Groundwater Transport (GWT) models with the Buoyancy (BUY) package so that density effects associated with the saltwater interface are represented. Flow in the combined sewer network is simulated with the Storm Water Management Model (SWMM) \citep{rossman2022swmm}.

The three simulators are advanced within a single process and exchange state through their respective library interfaces. MODFLOW~6 is controlled through its application programming interface, which extends the Basic Model Interface \citep{hughes2022modflowapi}, and D-Flow FM through its Basic Model Interface implementation \citep{hutton2020bmi}. SWMM implements neither, and is controlled through the SWMM toolkit by way of PySWMM \citep{mcdonnell2020pyswmm}; that interface opens a simulation from a SWMM input file, advances it by a routing step, closes it, and reads and writes node and link state, which are the same capabilities the other two supply under Basic Model Interface names. The MODFLOW~6 model is constructed with FloPy \citep{hughes2023flopy}. All three simulators are used as distributed, without modification.

% VERSION STATEMENT -- must match whichever runs produce the reported results.
% The completed 21-scenario sweep used MODFLOW 6 version 6.5.0; 6.7.0 is what is
% configured now. Reconcile before submission.
The versions used are D-Flow FM kernel 1.2.184, MODFLOW~6 version 6.7.0, and SWMM version 5.2.4 accessed through PySWMM 2.0.1.

\section{Coupling framework}

The coupled simulation is driven by the D-Flow FM model time. D-Flow FM advances in increments of its user time step, $\Delta t_u$, and MODFLOW~6 is advanced at a coupling interval, $\Delta t_c$, that is an integer multiple of $\Delta t_u$. The order of execution and the state passed between the models are shown in Figure~\ref{fig:sequence}.

Two cadences operate. On every user time step SWMM is advanced first, so that its outfall covers the interval D-Flow FM is about to integrate, and the mean discharge over that step is written into the D-Flow FM source term before D-Flow FM advances. Once per coupling interval, after the last user time step of the interval, the D-Flow FM water levels and depths are mapped to the MODFLOW~6 coastal boundary, MODFLOW~6 is advanced one time step, and the simulated boundary flows are returned to D-Flow FM. Figure~\ref{fig:sequence} is read by where each connector begins. A connector leaves the producing model time step at the point where the exchanged value is formed---the end of that step for a quantity taken instantaneously, the middle for a mean over one step, the whole span for a mean over an interval---and enters the start of the step that consumes it. Its horizontal run is therefore the offset between production and use, measured on the time axis.

The two directions of this exchange carry different lags. The water levels are averaged over the coupling interval, by equation~\ref{eq:tmean} below, and are the boundary condition for the MODFLOW~6 step spanning that same interval, so they introduce no lag. The boundary flows, in contrast, can only be recovered after MODFLOW~6 has finalized that step, so they are held constant over the \emph{following} coupling interval and lag the hydrodynamics by $\Delta t_c$. That lag is the principal coupling error, and it is what the sweep over coupling intervals reported below is designed to bound. A single MODFLOW~6 step therefore takes its coastal boundary from its own span and its sewer flux from the span before it. The aquifer--sewer seepage is returned on the same schedule as the boundary flows, and so carries the same lag. It is two-way: it is computed from the simulated heads and the water surface in the pipe together, and the same flux is applied to both models with opposite sign, as an inflow to the SWMM junction and an equal withdrawal from the MODFLOW~6 cell, so that mass leaving one model enters the other exactly. Because each exchange uses the state at the end of the preceding operation, the coupling is fully explicit and no iteration is performed within a coupling interval.

\begin{figure}[ht]
\centering
\widefig{figures/coupling_sequence.pdf}
\caption{Model time steps, order of execution, and state exchanged among the three simulators over two coupling intervals. Bar width is the model time step each simulator integrates: MODFLOW~6 advances the whole coupling interval, $\Delta t_c$, in one step, while SWMM and D-Flow FM tile that span with one bar per user time step, $\Delta t_u$. Lanes are ordered by execution rather than by physical position. Each connector leaves the producing step at the point where the exchanged value is formed and enters the start of the step that consumes it. Repeated user time steps are drawn faded, connector included, and only the leading one is labelled. The exchange break carries no model time---both of its edges have the same tick---and the faded stub at the right is the interval that follows. $Q_j$ is drawn in a neutral color because it is a property of neither model.}
\label{fig:sequence}
\end{figure}

The explicit scheme is imposed by what the interfaces expose, not chosen for convenience. An implicit coupling requires every participant to re-evaluate a time step against revised exchange terms, and only one of the three permits it: the MODFLOW application programming interface exposes the nonlinear solution within a time step, so the groundwater solution can be iterated \citep{hughes2022modflowapi}. The hydrodynamic and sewer interfaces do not. Both publish entry points that bracket a time step and advance it, with no facility to re-solve the same step once an exchanged quantity has changed, so a coupled iteration has nothing to iterate on two of its three legs.

The disparity in time step compounds the difficulty. The hydrodynamic model advances on a user time step of seconds to minutes, the groundwater model on the coupling interval, and the sewer model routes internally on a step of its own, so the models step at rates differing by up to two orders of magnitude. Even were re-evaluation available throughout, defining a converged coupled state across those scales, and deciding which model repeats how many of its own steps per outer iteration, is not a matter of simply looping the exchange. Quantifying how the explicit solution depends on the exchange interval, which the demonstration below does, is the useful thing that can be established without it.

The coupling interval is the principal numerical control on the approach. The coastal boundary of the groundwater model is forced by a semidiurnal tide with a period of approximately \SI{12.42}{\hour}, so a coupling interval that is a substantial fraction of that period samples the tidal signal sparsely and aliases it. Coupling intervals of 15 minutes, 30 minutes, and 1, 2, 4, 8, and 24 hours were simulated on each of three D-Flow FM grids to bound this effect.

\section{Coastal exchange}

Water levels simulated by D-Flow FM are applied to the MODFLOW~6 coastal boundary, which is represented by a General-Head Boundary (GHB) package covering the inner bays and a Constant-Head (CHD) package covering the open coast and the lateral perimeter. The D-Flow FM grid and the MODFLOW~6 grid are not coincident, so the exchange is mediated by area-weighted mapping matrices computed once, before simulation, from the intersection of the two grids.

D-Flow FM resolves the tide on a time step orders of magnitude shorter than the coupling interval, so a single boundary value has to stand for everything that happened over the interval. MODFLOW~6 advances that interval in one backward-in-time step, and the exchange term of such a step is applied across its whole length: the volume exchanged is the product of the interval, the conductance, and the head difference. The boundary value must therefore represent the interval rather than any instant within it, and the reduction used here is a time average weighted by the wetted state,

\begin{equation}
  \bar{\delta}_{j} = \frac{1}{\Delta t_{c}}
      \int_{t_{k}}^{t_{k+1}} \delta_{j}(t) \, \mathrm{d}t ,
  \qquad
  \bar{s}_{j} = \frac{\int_{t_{k}}^{t_{k+1}} \delta_{j}(t) \, s_{j}(t) \,
      \mathrm{d}t}{\int_{t_{k}}^{t_{k+1}} \delta_{j}(t) \, \mathrm{d}t} ,
  \label{eq:tmean}
\end{equation}

\noindent where $\bar{\delta}_{j}$ is the fraction of the coupling interval over which D-Flow FM cell $j$ is wet (dimensionless), $\delta_{j}(t)$ is unity where that cell is wet and zero where it is dry (dimensionless), $s_{j}(t)$ is the water level it simulates (L), $\bar{s}_{j}$ is the mean of that water level over the wetted part of the interval (L), and $\Delta t_{c}$ is the coupling interval (T). The integrals are evaluated by the trapezoidal rule on the D-Flow FM user time step. Weighting the level by the wetted state is necessary rather than tidy: an unweighted mean would credit an intertidal cell with a water level over the part of the interval it stood dry.

The boundary head assigned to GHB cell $i$ is

\begin{equation}
  h_{i}^{\mathrm{GHB}} = \gamma \sum_{j} w_{ij}^{\mathrm{GHB}} \, \bar{s}_{j} ,
  \label{eq:ghbhead}
\end{equation}

\noindent where $h_{i}^{\mathrm{GHB}}$ is the boundary head of cell $i$ (L), $w_{ij}^{\mathrm{GHB}}$ is the area-weighted contribution of D-Flow FM cell $j$ to MODFLOW~6 cell $i$ (dimensionless), and $\gamma$ is the unit conversion from meters to feet (dimensionless). Constant-head values are computed identically using the corresponding CHD mapping.

Intertidal cells alternately wet and dry over a tidal cycle. A D-Flow FM cell with zero water depth contributes no head and no conductance, so the GHB conductance is reduced in proportion to the wetted fraction of the contributing cells,

\begin{equation}
  C_{i}^{\mathrm{GHB}} = C_{i}^{0} \sum_{j} w_{ij}^{\mathrm{GHB}} \,
      \bar{\delta}_{j} ,
  \label{eq:ghbcond}
\end{equation}

\noindent where $C_{i}^{\mathrm{GHB}}$ is the conductance applied to cell $i$ (L$^{2}$T$^{-1}$), and $C_{i}^{0}$ is the conductance of cell $i$ in the base model (L$^{2}$T$^{-1}$). Equations~\ref{eq:tmean} to \ref{eq:ghbcond} together reproduce the interval-mean exchange exactly rather than approximately. Because the instantaneous conductance is the base conductance multiplied by a quantity that is either zero or one, and because MODFLOW~6 carries a single head over the step,

\begin{equation}
  \frac{1}{\Delta t_{c}} \int_{t_{k}}^{t_{k+1}} C^{0} \delta_{j}(t)
      \left( s_{j}(t) - h \right) \mathrm{d}t
    = C^{0} \, \bar{\delta}_{j} \left( \bar{s}_{j} - h \right) ,
  \label{eq:exact}
\end{equation}

\noindent where $h$ is the groundwater head at the end of the step (L). The wetted fraction is thus the conductance multiplier and the wet-weighted mean is the boundary head, and no other pair of values reproduces the interval-mean flux. The base-model conductance, rather than the conductance from the preceding coupling interval, is used on the right-hand side of equation~\ref{eq:ghbcond}; applying the reduction multiplicatively to the previous value would drive the conductance monotonically toward zero over the simulation.

Boundary flows simulated by MODFLOW~6 are returned to D-Flow FM as an external volume source. The source applied to D-Flow FM cell $j$ is

\begin{equation}
  Q_{j}^{\mathrm{ext}} = -\beta \, \delta_{j}
    \left( \sum_{i} a_{ji}^{\mathrm{GHB}} Q_{i}^{\mathrm{GHB}}
         + \sum_{i} a_{ji}^{\mathrm{CHD}} Q_{i}^{\mathrm{CHD}} \right) ,
  \label{eq:qext}
\end{equation}

\noindent where $Q_{j}^{\mathrm{ext}}$ is the volumetric source applied to D-Flow FM cell $j$ (L$^{3}$T$^{-1}$), $a_{ji}$ is the area-weighted contribution of MODFLOW~6 boundary $i$ to D-Flow FM cell $j$ (dimensionless), $Q_{i}$ is the boundary flow simulated by MODFLOW~6 (L$^{3}$T$^{-1}$), and $\beta$ is the unit conversion from cubic feet per day to cubic meters per second (dimensionless). The sign convention of MODFLOW~6 reports boundary flow as positive into the groundwater model, so the leading negative sign renders $Q_{j}^{\mathrm{ext}}$ positive into the surface water. Dry cells receive no source.

\subsection{Choice of temporal reduction}
\label{sec:reduction}

Averaging is one of two reductions in common use. The other is to sample the D-Flow FM state at the end of the coupling interval and apply it unchanged, which is the simpler implementation and is consistent in the sense that it converges as the interval is shortened. The two are not better and worse versions of one approximation. They fail by different mechanisms, and which mechanism is active is set by the tide rather than by the model. Sampling preserves the amplitude of the boundary signal but aliases it once the coupling interval exceeds the Nyquist limit of the dominant tidal constituent. Averaging cannot alias, but damps the amplitude by an amount that grows with the interval.

Figure~\ref{fig:reduction} shows what each reduction does to cells of differing wetness. The water level used there is the sum of the three semidiurnal constituents fitted by least squares to the National Oceanic and Atmospheric Administration record at Kings Point, New York (station 8516945, referenced to NAVD 1988), in mid Long Island Sound: \SI{1.110}{\meter} for $M_2$, \SI{0.225}{\meter} for $S_2$, and \SI{0.287}{\meter} for $N_2$. The spring--neap and perigean beat are therefore those of the measured record rather than of a single sinusoid, which matters because it is the beat between that record and the coupling interval that produces the behavior described below. For a cell that is never dry, the conductance is fixed and only the head is at issue, and the trade between aliasing and damping is the whole of it. For an intertidal cell the conductance carries the more consequential difference. Sampled instantaneously, such a cell is reported as either fully connected or entirely disconnected depending on where the sampling instant happens to fall in the tide, so the boundary switches on and off at the beat frequency between the coupling interval and the tide while nothing physical about the cell has changed. The wetted fraction of equation~\ref{eq:tmean} instead reports the partial connection the interval actually contained.

\begin{figure}[ht]
\centering
\widefig{figures/boundary_reduction.pdf}
\caption{Reduction of a tidal boundary to one value per coupling interval, for cells wet 100, 75, 50, and 25 percent of the time, at coupling intervals of 2, 8, and 24 hours. Each pair of panels shows the water level and, beneath it, the multiplier applied to the base conductance. Bed elevations are placed at quantiles of the water level to give the stated wetted fractions. Crosses mark coupling intervals for which instantaneous sampling reports the cell dry, so that no boundary is applied at all. The 2-hour interval is below the $M_2$ Nyquist limit of \SI{6.21}{\hour} and the other two are above it.}
\label{fig:reduction}
\end{figure}

The consequence for the coupled solution was evaluated by simulating the coarse grid at coupling intervals of 30 minutes and 1, 2, 4, 6, 8, 12, and 24 hours under both reductions, and scoring each against a 15-minute simulation, which resolves the tide and serves as the reference. Every interval divides a day, so that coupling steps coincide with the boundaries of the daily stress periods. The medium and fine grids were simulated over the same set excluding the 6- and 12-hour intervals, which places the principal result on three independent discretizations. Both a sampled and an averaged 15-minute simulation were used as the reference, because the averaged one shares its formulation with half of the simulations being scored; the two give the same result, so the comparison does not depend on that choice.

Results are summarized in Figure~\ref{fig:averaging}. At daily coupling, where the choice of reduction has consequence, sampling carries \SI{33}{\milli\meter} of head error against \SI{3}{\milli\meter} for averaging on the coarse grid. The medium and fine grids give \SI{39}{\milli\meter} against \SI{3}{\milli\meter} and \SI{31}{\milli\meter} against \SI{3}{\milli\meter}, ratios of 12.0, 11.4, and 10.4, so the result does not depend on the discretization. At intervals of 8 hours and shorter the two reductions differ by amounts too small to matter---root-mean-square head error stays below \SI{2.2}{\milli\meter} on all three grids---and there sampling is the marginally more accurate of the two at fifteen of the sixteen combinations of grid and interval, which is the damping penalty that averaging pays. Aquifer--sewer seepage behaves as the head does, favoring averaging at daily coupling by factors of 1.4, 2.9, and 2.7. The sewer discharge tracer does not distinguish the two reductions at all: its 99.99th percentile stays within 1 percent of the reference at every interval on the coarse grid and within a third of a percent on the medium grid, and on the fine grid sits 3 to 4 percent from it at intervals where the head has already converged, reaching 18 percent only at daily coupling.

What distinguishes the two error mechanisms is that only one of them is monotonic in the coupling interval. A truncation error grows smoothly as the step is lengthened, and the averaged solution does exactly that, from \SI{0.05}{\milli\meter} at 30 minutes to \SI{2.7}{\milli\meter} at daily coupling. The sampled solution does not: its head error is \SI{1.3}{\milli\meter} at 8 hours, \SI{36.9}{\milli\meter} at 12 hours, and \SI{32.9}{\milli\meter} at 24 hours, so halving the coupling interval from one day to twelve hours makes the solution worse rather than better. Sampling every 12 and every 24 hours folds $M_2$ onto the same 14.8-day beat, and the two carry comparable error; sampling every 8 hours folds it onto a 22.5-hour beat, which the aquifer damps and which costs almost nothing even though 8 hours is already above the Nyquist limit. It is the period onto which the constituent is aliased, rather than the length of the interval or the crossing of the limit in itself, that sets the error. Because a coupling interval is normally chosen for computational cost rather than for its relation to the tide, and because the penalty cannot be relied upon to fall as the interval is shortened, averaging is used throughout this work.

\begin{figure}[ht]
\centering
\widefig{figures/boundary_averaging.pdf}
\caption{Error in the coupled solution relative to a 15-minute reference simulation, for the two reductions of the D-Flow FM boundary, on the coarse grid. (A) simulated groundwater head, (B) aquifer--sewer seepage, (C) sewer discharge tracer, and (D) the 99.99th percentile of the sewer tracer, as a percent departure from the reference shown dashed. The first five days are excluded as model spin-up. The vertical rule is the Nyquist limit for the $M_2$ constituent, \SI{6.21}{\hour}. The sampled head error in (A) is not monotonic in the coupling interval: it is larger at 12 hours than at 24, because both intervals alias $M_2$ onto the same 14.8-day beat, whereas the 8-hour interval aliases it onto a 22.5-hour beat that the aquifer damps. The averaged error grows smoothly throughout, as a truncation error does. Panel (D) is an upper-tail percentile rather than the maximum of the field, which falls in a single cell and is not a domain-wide measure.}
\label{fig:averaging}
\end{figure}

\section{Sewer exchange}

Exchange between the aquifer and the sewer is computed at a set of junctions that fall within the active area of the groundwater model. Each junction is associated with the MODFLOW~6 cell containing it, and a well is added to the groundwater model at that cell.

Infiltration and exfiltration are driven by the difference between the groundwater head and the water surface inside the pipe, and are limited by a conductance representing the leakage area of the pipe wall served by the junction. A per-junction conductance is computed from the sewer geometry as

\begin{equation}
  C_{j} = \lambda \, \pi D_{j} L_{j} ,
  \label{eq:pipecond}
\end{equation}

\noindent where $C_{j}$ is the leakage conductance of junction $j$ (L$^{2}$T$^{-1}$), $\lambda$ is a bulk pipe-wall leakance (T$^{-1}$), $D_{j}$ is the conduit diameter (L), and $L_{j}$ is the length of pipe served by junction $j$ (L), taken as half the length of every conduit that the junction touches. The served length varies by more than a factor of twenty across the network, from approximately 28 to \SI{677}{\foot}, so a conductance assigned uniformly per connection cannot represent the network geometry.

The driving head depends on whether the water table reaches the pipe. Where the groundwater head stands above the pipe invert the surrounding material is saturated and the gradient is measured to the water surface within the pipe; where the head falls below the invert the pipe free-drains through unsaturated material and the gradient is set by the depth of water in the pipe alone. The driving head is therefore

\begin{equation}
  \Delta h_{j} =
  \begin{cases}
    h_{j} - \left( z_{j} + d_{j} \right), & h_{j} > z_{j} \\[4pt]
    -d_{j},                                & h_{j} \le z_{j}
  \end{cases}
  \label{eq:pot}
\end{equation}

\noindent where $\Delta h_{j}$ is the driving head at junction $j$ (L), $h_{j}$ is the simulated groundwater head in the cell containing junction $j$ (L), $z_{j}$ is the pipe invert elevation (L), and $d_{j}$ is the depth of water in the pipe above the invert (L). Positive $\Delta h_{j}$ produces infiltration into the sewer and negative $\Delta h_{j}$ produces exfiltration. Treating the disconnected case as $-d_{j}$ rather than $h_{j} - z_{j}$ is the same treatment given a disconnected stream, and it is necessary: the latter form produces exfiltration from an empty pipe that grows without bound as the water table falls.

The exchange at junction $j$ is

\begin{equation}
  Q_{j} = \eta_{j} C_{j} \Delta h_{j} ,
  \label{eq:swmmq}
\end{equation}

\noindent where $Q_{j}$ is the exchange (L$^{3}$T$^{-1}$), and $\eta_{j}$ is unity for infiltration and an exfiltration ratio less than unity for exfiltration (dimensionless). Exfiltration into saturated material is more difficult than infiltration from it, and the asymmetry is retained. The exchange is applied to SWMM as a generated inflow at the junction and to MODFLOW~6 as an equal and opposite well flux, so that mass is conserved between the two models. Cells that are dry or inactive return a head of $-10^{30}$ in MODFLOW~6 and are assigned zero exchange rather than being used unguarded.

\section{Sewer discharge tracer transport}

Flow reaching the outfall is introduced to D-Flow FM as a point source carrying a conservative tracer of sewer discharge, which is then advected and dispersed by the hydrodynamic model. The source discharge is the total inflow simulated by SWMM at the outfall, and the tracer concentration assigned to that discharge is constant, so the tracer mass flux follows the simulated SWMM outfall discharge directly.

The SWMM outfall discharge is written into the D-Flow FM source term at each user time step. SWMM is advanced one user time step ahead of D-Flow FM so that its outfall covers the interval D-Flow FM is about to integrate, and the discharge applied over that interval is the mean

\begin{equation}
  \bar{Q}_{s} = \frac{V\!\left(t + \Delta t_{u}\right) - V\!\left(t\right)}
                     {\Delta t_{u}} ,
  \label{eq:blockmean}
\end{equation}

\noindent where $\bar{Q}_{s}$ is the mean SWMM outfall discharge over the user time step (L$^{3}$T$^{-1}$), $V$ is the cumulative volume conveyed through the SWMM outfall (L$^{3}$), and $\Delta t_{u}$ is the D-Flow FM user time step (T). Forming the mean from the cumulative volume conserves mass exactly over the step, whereas sampling the SWMM outfall discharge instantaneously once per coupling interval and holding it would introduce the same aliasing that the coupling interval itself is being tested for. Advancing SWMM at the user time step rather than at the coupling interval does not alter the sewer solution, because SWMM routes internally at its own routing step in either case and the exchange computed from equation~\ref{eq:swmmq} is still refreshed only once per coupling interval; it raises the resolution at which the outfall is observed.

\section{Implementation}

The coupled simulation is implemented in Python. The three simulators are loaded as shared libraries into a single process, so that state is passed between them as arrays in memory.

D-Flow FM applies its external forcings at the start of a user time step and then computes the flow solution. A source term written from outside the model between those two operations is used by the solution, whereas one written around the complete time step is overwritten by the forcing update before it is reached. The direct tracer mechanism therefore advances D-Flow FM through the separate initialization, computation, and finalization entry points of its interface rather than through the combined update, and writes the source term between the first and the second. This is the transfer drawn as $\bar{Q}_s$ in figure~\ref{fig:sequence}; the arrow falls within a user time step, not between two of them.

This is the sub-time-step access described earlier, and the two simulators reach it differently. MODFLOW~6 extends the Basic Model Interface with functions that expose the time step and the solution within it \citep{hughes2022modflowapi}, and those functions are what the groundwater exchange is written through. D-Flow FM instead publishes its user time step as three separate entry points alongside the combined one. Neither route is part of the Basic Model Interface itself, so a coupling of this kind is portable only to the extent that each simulator provides some equivalent.

\subsection{Coupling strategy as a property of the driver}

Because no simulator is modified, the coupling strategy---what is exchanged, how it is reduced to a single value, and where within a model time step it is written---resides entirely in the driver. An alternative strategy is then a change to a few lines of Python and a re-run, rather than an edit and a rebuild of a hydrodynamic or a groundwater code. Two of the choices reported here were settled that way, and neither would have been convenient to settle otherwise.

The reduction of the D-Flow FM boundary to one value per coupling interval (Section~\ref{sec:reduction}) is selected by a single setting, so the instantaneous and the time-averaged forms could both be simulated, at every coupling interval, using the same executables and the same model input. That is what makes the comparison in Figure~\ref{fig:averaging} a controlled one: a single thing differs between the two curves, and the aliasing threshold it exposes is a property of the reduction rather than of the build. Implementing either formulation inside a simulator would have required a separate build for each, and would have left the two unavailable simultaneously.

Where in the D-Flow FM time step the sewer source term has to be written was established the same way. Writing it around the combined time-stepping call delivered none of the outfall volume, because the external forcing update overwrites the source term before the flow computation reads it; writing it between the separate initialization and computation entry points delivered all of it. That is an experiment performed on the driver, against the distributed library, requiring no knowledge of the internals beyond the names of the entry points and no ability to rebuild the code.

The freedom is bounded by what the interfaces expose, as the preceding paragraphs describe: a strategy that needs state the simulator does not publish, or needs it at a moment the simulator does not expose, cannot be written in the driver at all. Within those bounds, however, the cost of trying a coupling strategy falls to the cost of running the models, which is what allowed the exchange formulations here to be chosen on evidence rather than on argument.

\section{Demonstration}

The coupling is exercised on the Greenport sewer network described above. The groundwater model is a cutout of the regional Long Island aquifer system model \citep{walter2024longisland}, which simulates the period 1900 to 2019 with MODFLOW~6 including the position of the saltwater interface. Three D-Flow FM grids of increasing resolution were used---6,491, 16,666, and 41,091 cells---spanning the same domain, sharing a user time step of \SI{300}{\second}, and driven by the same tidal and atmospheric forcing over the same 89-day period. The groundwater and sewer models are identical across all simulations, so differences among grids reflect the resolution of the coastal water level alone, and differences among exchange intervals reflect the coupling alone. Each grid was simulated at seven coupling intervals---15 and 30 minutes and 1, 2, 4, 8, and 24 hours---under both reductions of the coastal boundary, with two further intervals, 6 and 12 hours, added on the coarse grid, giving 46 simulations in total. Every interval divides a day, so that coupling steps coincide with the boundaries of the daily stress periods.




\section{Results}

The sensitivity of the coupled solution to the exchange interval is reported in Section~\ref{sec:reduction} and is not repeated here. Two of its features bear on the requirements a coupled model imposes, and both survive a change of grid. The first is that the penalty for sampling the coastal boundary rather than averaging it over the interval is a factor of 12.0, 11.4, and 10.4 in aquifer head at daily coupling on the coarse, medium, and fine grids, with the aquifer--sewer seepage agreeing at 1.4, 2.9, and 2.7; three independently discretized solutions place the same requirement on the coupling. The second is that the sewer discharge tracer does not resolve the difference at all. Its upper-tail concentration stays within 1 percent of the reference at every interval on the coarse grid and within a third of a percent on the medium grid, and departs materially only at daily coupling on the fine grid. It is the aquifer response, not the transported quantity, that the exchange interval is felt in.

The requirement the coupling places on the exchange interval does not depend on the discretization, but the exchange itself does. Figure~\ref{fig:sewer} compares the aquifer--sewer exchange simulated by the three grids at a 15-minute interval, where the coupling contributes least and what remains is hydrodynamic resolution. Net infiltration over the simulation is 180, 79, and 87 gallons per day per inch of pipe diameter per mile of sewer on the coarse, medium, and fine grids. The coarse grid returns more than twice the medium one, and the sequence is not monotonic: refining from medium to fine raises the rate by 10 percent rather than continuing to lower it. Against an acceptance range of 50 to 200 in those units for new gravity sanitary sewers, the choice of hydrodynamic grid moves the simulated rate across much of the band.

The mechanism is in Figure~\ref{fig:sewer}B. What differs between the grids is the fraction of junctions standing below the water table, and so exchanging at all, which averages 23.6, 18.0, and 18.7 percent over the same period and follows the same non-monotonic ordering. Coastal water levels set the head along the seaward boundary, the head sets which junctions are connected, and the connected fraction sets the exchange. The sensitivity is amplified along that path rather than merely transmitted: a 24 percent reduction in connected junctions between the coarse and medium grids produces a 56 percent reduction in net infiltration.

Two consequences follow. The sewer exchange is not converged in hydrodynamic resolution over the range simulated here, so an infiltration rate obtained from a coupled model of this kind should be reported together with the grid that produced it. And the bulk pipe-wall leakance is not transferable between grids: it is held at \SI{0.0655}{\per\day} in all three simulations so that the differences are resolution alone, with the consequence that a leakance calibrated to reproduce an observed rate on one grid would be wrong by approximately a factor of two on another.

What the exchange interval costs is shown in Figure~\ref{fig:cost}. The marginal cost of one coupling step is \SI{0.25}{\second} and does not depend on the grid, which is what should be expected: the MODFLOW~6 and SWMM models are identical across the three D-Flow FM discretizations, so a shorter interval buys additional solves of the same two models. The hydrodynamic cost it is set against is not constant. D-Flow FM alone accounts for approximately 36, 94, and \SI{239}{\minute} of the 89-day simulation on the coarse, medium, and fine grids, which is very nearly linear in cell count at 5.5 to $5.9\times10^{-3}$ minutes per cell. The consequence is that the penalty for frequent coupling falls as the hydrodynamic model is refined: increasing the number of exchanges ninety-six-fold, from daily to 15-minute coupling, costs of order 100 percent of the run on the coarse grid but between 13 and 39 percent on the medium and fine grids.

That asymmetry is the practical result of this work. The error incurred by a long exchange interval is not monotonic in the interval, so it cannot be bounded by shortening the interval a little; the cost of a short one is modest, and becomes more modest as the hydrodynamic grid approaches the resolution at which such models are actually run. There is correspondingly little to recommend a long exchange interval on a production-scale grid, and the coupling interval is better chosen from the tidal constituents than from the run-time budget.

Run times were measured on a single workstation with the simulations executed sequentially. Run-to-run variation is a few percent, which is the reason several medium-grid points in Figure~\ref{fig:cost}A fall marginally below their own daily-coupling baseline, and the reason the spread between the two reductions at a given grid should be read as scatter rather than as a property of the reduction.

\begin{figure}[ht]
\centering
\widefig{figures/sewer_exchange.pdf}
\caption{Aquifer--sewer exchange simulated by the three D-Flow FM grids at a 15-minute coupling interval under the averaged coastal boundary. (A) cumulative net exchange into the sewer, annotated with the infiltration rate each grid settles at, in gallons per day per inch of pipe diameter per mile; positive is infiltration. (B) percentage of the 244 coupled junctions standing below the water table and therefore exchanging, plotted at every coupling step so that the tidal swing is visible. The first five days, shaded, are excluded from the reported rates as model spin-up. The MODFLOW~6 and SWMM models are identical in the three simulations and the bulk pipe-wall leakance is the same, so the differences are those of hydrodynamic resolution alone.}
\label{fig:sewer}
\end{figure}

\begin{figure}[ht]
\centering
\widefig{figures/coupling_cost.pdf}
\caption{Cost of the exchange interval, for both reductions of the D-Flow FM boundary on the three grids. (A) run time relative to that series' own daily-coupling simulation. (B) additional run time against additional coupling steps, each series referred to its own daily-coupling simulation rather than to a fitted intercept, so that the collapse shown is a property of the measurements; the line is the pooled marginal cost through the origin. The vertical rule in (A) is the Nyquist limit for the $M_2$ constituent, \SI{6.21}{\hour}. Coupling more frequently costs approximately the same number of seconds per exchange on every grid, because the MODFLOW~6 and SWMM models being solved are identical; what differs is the hydrodynamic cost that expense is measured against.}
\label{fig:cost}
\end{figure}

\section{Summary and conclusions}

% Placeholder. Prose paragraphs, not bullets, once the analysis is complete:
% the coupling approach and what it costs; the character and bounds of the error
% introduced by the exchange interval, and how frequently the exchange must occur;
% and what the demonstration establishes about the requirements a library
% interface must meet for this kind of coupling -- that sub-time-step access is
% needed, that the Basic Model Interface does not define it, and that both
% simulators here supply it by different and non-portable routes.
%
% Forward-looking paragraphs close the section:
%   - what an implicit coupling would require that is not available today: the
%     hydrodynamic and sewer interfaces advance a time step but cannot re-solve
%     one, so two of the three legs have nothing to iterate on, and the models
%     step at rates differing by up to two orders of magnitude. This is a request
%     of the simulators, not a limitation of the approach.
%   - applying the same pattern to other simulator combinations, which requires
%     only that each expose an equivalent sub-time-step entry point

% Journal-specific and expected of most EMS papers: name of software or dataset,
% developer and contact, year first available, hardware and software required,
% availability and cost, program language and size for software, and repository
% form, archive size and access form for data. "Contact the author" is explicitly
% not acceptable. Sizes, the coupling-code archive DOI, and the data deposit are
% still to be supplied.
\section*{Software and data availability}

The three simulators are used as distributed and are not modified by this work.
MODFLOW~6 is developed by the U.S. Geological Survey, is written in Fortran, and is
available at no cost from \url{https://github.com/MODFLOW-ORG/modflow6}; version
6.7.0 was used, with the API and the \texttt{modflowapi} Python interface. D-Flow FM
is developed by Deltares, is written in Fortran, and is available from
\url{https://oss.deltares.nl/web/delft3dfm}. SWMM is developed by the
U.S. Environmental Protection Agency, is written in C, and is available at no cost
from \url{https://github.com/USEPA/Stormwater-Management-Model}; version 5.2.4 was
used, driven through the \texttt{PySWMM} Python interface, version 2.0.1. All three
are used as shared libraries on 64-bit Windows and Linux; no special hardware is
required.

The coupling itself is Python and is the only software new to this work. % TODO: archive the coupling code and cite the DOI here rather than a bare repository URL.

% TODO: EMS is Option C for research data -- the model input and the simulation
% output must be deposited in a repository and cited, or a statement given
% explaining why they cannot be shared. Neither is done yet.

\section*{CRediT authorship contribution statement}

% TODO: required by the guide, and only the co-authors can supply it. Roles are
% drawn from the CRediT taxonomy: conceptualization, data curation, formal analysis,
% funding acquisition, investigation, methodology, project administration,
% resources, software, supervision, validation, visualization, writing -- original
% draft, writing -- review and editing.

\section*{Declaration of generative AI and AI-assisted technologies in the manuscript preparation process}

% TODO: an author decision, not a drafting one. The guide requires this section
% when generative AI tools were used in preparing the manuscript, placed before the
% reference list, and it is omitted entirely only if there is nothing to disclose.
% Elsevier's wording is: "During the preparation of this work the author(s) used
% [NAME OF TOOL] in order to [REASON]. After using this tool/service, the author(s)
% reviewed and edited the content as needed and take(s) full responsibility for the
% content of the published article." The guide separately notes that AI tools may be
% used for conceptual diagrams and for data visualizations derived from underlying
% data by reproducible methods, and that such use should be disclosed in the image
% caption as well as here.

\section*{Acknowledgments}

% Placeholder.

\section*{Authors' Note}

The authors do not have any conflicts of interest to report.

\section*{Disclaimer}

Any use of trade, firm, or product names is for descriptive purposes only and does not imply endorsement by the U.S. Government.

\bibliographystyle{elsarticle-harv}
\bibliography{Hughesetal_ESM_LISSCoupling}

\end{document}
